\documentclass[10pt]{article}
\usepackage[T1]{fontenc}
\usepackage[utf8]{inputenc}
\usepackage[french]{babel}
\usepackage{amsmath,amssymb,mathrsfs,mathpazo}
\usepackage{geometry}
\geometry{a4paper,margin=2cm}
\usepackage{graphicx,array,tabularx,tabularray,colortbl,multirow,multicol,booktabs,enumitem,pifont,ulem,fancybox,diagbox,cancel}
\usepackage{pstricks,pstricks-add,pst-plot,pst-tree,pst-node,pst-text,pst-func,pst-math,pst-bspline,pst-3dplot,pst-eucl,pst-solides3d}
\let\pscalculate\undefined
\usepackage{pst-bezier}
\usepackage{tikz,pgfplots,tkz-tab}
\usetikzlibrary{arrows,arrows.meta,patterns,calc,positioning,trees,decorations.markings,shapes.geometric,shapes.misc,shapes.arrows,intersections,decorations.pathmorphing,decorations.pathreplacing,backgrounds,shadows,babel}
\pgfplotsset{compat=1.18}
\usepgfplotslibrary{fillbetween}
\usepackage[np]{numprint}
\usepackage{listings,makecell,pas-tableur,twoopt}
\let\barre\undefined
\usepackage[pstricks]{tablvar}
\definecolor{darkgreen}{RGB}{0,100,0}
\providecommand{\red}{\color{red}}
\providecommand{\blue}{\color{blue}}
\providecommand{\green}{\color{green}}
\providecommand{\black}{\color{black}}
\providecommand{\white}{\color{white}}
\providecommand{\gray}{\color{gray}}
\providecommand{\cyan}{\color{cyan}}
\providecommand{\magenta}{\color{magenta}}
\providecommand{\yellow}{\color{yellow}}
\providecommand{\violet}{\color{violet}}
\providecommand{\orange}{\color{orange}}
\providecommand{\euro}{\texteuro}
\providecommand{\EUR}{\texteuro}
\providecommand{\ieme}{\textsuperscript{e}}
\providecommand{\degres}{\ensuremath{^\circ}}
\providecommand{\degree}{\ensuremath{^\circ}}
\providecommand{\vv}[1]{\overrightarrow{#1}}
\providecommand{\Coord}[1]{#1}
\providecommand{\decosix}{}
\providecommand{\starredbullet}{\ensuremath{\bullet}}
\pagestyle{empty}
\setlength{\parindent}{0pt}
\usepackage{adjustbox}\DeclareUnicodeCharacter{2236}{:}\DeclareUnicodeCharacter{2190}{\ensuremath{\leftarrow}}
\providecommand{\R}{\mathbb{R}}
\providecommand{\N}{\mathbb{N}}
\providecommand{\Z}{\mathbb{Z}}
\providecommand{\Q}{\mathbb{Q}}
\providecommand{\D}{\mathbb{D}}
\providecommand{\C}{\mathbb{C}}
\providecommand{\e}{{\mathrm{e}}}
\providecommand{\vect}[1]{\overrightarrow{#1}}
\providecommand{\vectt}[1]{\overrightarrow{\mathrm{#1}}}
\providecommand{\vvt}[1]{\overrightarrow{\mathrm{#1}}}
\providecommand{\np}[2][]{{#2~#1}}
\providecommand{\vv}[1]{{\overrightarrow{#1}}}
\renewcommand{\d}{{\mathrm{d}}}
\providecommand{\degres}{{^\circ}}
\providecommand{\degre}{{^\circ}}
\providecommand{\blue}{\color{blue}}
\providecommand{\red}{\color{red}}
\providecommand{\enskip}{\quad}
\providecommand{\llbracket}{\lbrack\!\lbrack}
\providecommand{\rrbracket}{\rbrack\!\rbrack}
\providecommand{\euro}{{\text{€}}}
\providecommand{\no}{{\text{n°}}}
\providecommand{\interval}[2]{{[#1;#2]}}
\providecommand{\intervFO}[2]{{[#1;#2[}}
\providecommand{\dotfill}{\ldots}

\providecommand{\ds}{\displaystyle}
\providecommand{\textding}[1]{#1}
\providecommand{\Oij}{\(\mathrm{O};\vec{\imath},\vec{\jmath}\)}
\providecommand{\Oijk}{\(\mathrm{O};\vec{\imath},\vec{\jmath},\vec{k}\)}
\providecommand{\Ouv}{\(\mathrm{O};\vec{u},\vec{v}\)}
\providecommand{\intervFF}[2]{\ensuremath{[#1;#2]}}
\renewcommand{\euro}{\ifmmode\text{\texteuro}\else\texteuro\fi}
\DeclareUnicodeCharacter{00B0}{\ifmmode^{\circ}\else\textdegree\fi}
\DeclareUnicodeCharacter{20AC}{\ifmmode\text{\texteuro}\else\texteuro\fi}
% Énoncés : Éducation nationale. Transcription et corrigés : APMEP, Denis Vergès et contributeurs. Découpage et classement : CoopMaths, Sébastien Lozano, Éric Elter et contributeurs. Adaptation de la présentation : Maths Explorer.
% CC BY-SA : https://coopmaths.fr/www/about/
% https://coopmaths.fr/alea/static/e3c/2020/tex/e3c_2020_00_sujet43_1_cor.tex
\begin{document}
\noindent\textit{APMEP / CoopMaths. CC BY-SA. Présentation adaptée par Maths Explorer.}\par\medskip
\subsection*{Question 1}

Puisque \( x > 2{,}5 \), alors \( -2x + 5 \neq 0 \), donc \( f \) est dérivable sur \( ]2{,}5 \,;\, +\infty[ \) et on a :
\begin{align*}
f'(x) &= \dfrac{3(-2x + 5) - (-2)(3x + 1)}{(-2x + 5)^2} \\
&= \dfrac{-6x + 15 + 6x + 2}{(-2x + 5)^2} \\
&= \dfrac{17}{(-2x + 5)^2}.
\end{align*}

\subsection*{Question 2}

La courbe coupe l'axe des abscisses aux points d'abscisse 1 et 3. L'expression de \( f(x) \) contient donc les facteurs \( x - 1 \) (ou \( 1 - x \)) et \( x - 2 \) (ou \( 2 - x \)).

D'autre part, on doit avoir \( f(2) = 2 \) : seule la réponse \textbf{b.} convient.

\subsection*{Question 3}

Le nombre dérivé \( f'(0) \) est égal au coefficient directeur de la tangente au point d'abscisse 0. En utilisant les points \( (0\,;\,2) \) et \( (2\,;\,4) \), on trouve que :
\[
f'(0) = \dfrac{4 - 2}{2 - 0} = \dfrac{2}{2} = 1.
\]

\subsection*{Question 4}

Une équation réduite de la droite \( (GH) \) est : \( y = ax + b \).

On a :
\[
\begin{cases}
G \in (GH) \\
H \in (GH)
\end{cases}
\iff 
\begin{cases}
-2 = a + b \\
4 = 6a + b
\end{cases}
\Rightarrow 6 = 5a \iff a = \dfrac{6}{5}.
\]
On en déduit que : \( b = -2 - \dfrac{6}{5} = -\dfrac{16}{5} \).

Une équation de la droite \( (GH) \) est : \( y = \dfrac{6}{5}x - \dfrac{16}{5} \).

On a :
\[
D(-14 \,;\, -20) \in (GH) \iff -20 = \dfrac{6}{5} \times (-14) - \dfrac{16}{5},
\]
qui est vraie.


\subsection*{Question 5}

De la relation \( \sin^2 x + \cos^2 x = 1 \), on en déduit que :
\[
\sin^2 x = 1 - \left( -\dfrac{\sqrt{3}}{2} \right)^2 = 1 - \dfrac{3}{4} = \dfrac{1}{4}.
\]
Or, sur l'intervalle \( \left[ \pi \,;\, \dfrac{3\pi}{2} \right] \), on sait que \( \sin x < 0 \), donc \( \sin x = -\dfrac{1}{2} \).

\end{document}
